\begin{abstract}
A distribuição normal ocupa posição central em muitos cursos introdutórios de tratamento de dados, embora muitas delas sejam assimétricas ou limitadas. Este artigo apresenta uma abordagem didática para a modelagem probabilística de dados físicos reais com a distribuição lambda generalizada, na parametrização FKML, definida por uma função quantil explícita. Explica-se o significado dos quatro parâmetros, as condições de validade e de suporte e o ajuste pelo método dos momentos, implementado na biblioteca Python \texttt{pyglam}. O procedimento é aplicado a dois conjuntos de dados de uma estação meteorológica automática do INMET em Brasília, que são a velocidade média diária do vento, comparada às distribuições de Rayleigh e de Weibull, e a temperatura média diária de julho, comparada à distribuição normal. A avaliação usa a distância de Kolmogorov-Smirnov e os desvios de quantis nas unidades da grandeza, interpretados à luz da resolução das medições e da incerteza amostral. No vento, a distribuição flexível tem os menores desvios nos quantis examinados, da mesma ordem da variabilidade estimada por reamostragem, mas seu suporte exclui 0,95\% dos dias, justamente os de vento mais fraco. Na temperatura, as discrepâncias da normal são menores que a variabilidade amostral dos quantis, e o modelo de dois parâmetros continua sendo uma escolha razoável para essa descrição. Os exemplos mostram que um bom ajuste exige examinar o suporte e a aderência, e não apenas os momentos, e que a escolha do modelo depende da pergunta física. Todos os resultados são reproduzidos por um notebook disponibilizado com o artigo, e uma sugestão de roteiro para uso em sala, com tarefas e respostas esperadas, acompanha o texto.

\medskip
\noindent\textbf{Palavras-chave}: distribuição lambda generalizada; função quantil; método dos momentos; dados meteorológicos; ensino de estatística; Python.
\end{abstract}

\begin{otherlanguage}{english}
\begin{abstract}
\noindent\textbf{Flexible modelling of physical phenomena with the generalized lambda distribution: a didactic computational approach with pyGLAM.}
The normal distribution has a central place in many introductory data-analysis courses, although many of them are skewed or bounded. This paper presents a didactic approach to the probabilistic modelling of real physical data with the generalized lambda distribution, in the FKML parameterization, defined by an explicit quantile function. It explains the meaning of the four parameters, the conditions of validity and support, and the fit by the method of moments, implemented in the Python library \texttt{pyglam}. The procedure is applied to two data sets from an automatic weather station of INMET in Brasília, Brazil, namely the daily mean wind speed, compared with the Rayleigh and Weibull distributions, and the daily mean temperature in July, compared with the normal distribution. The assessment uses the Kolmogorov-Smirnov distance and quantile deviations in the units of the quantity, interpreted in the light of the measurement resolution and of the sampling uncertainty. For the wind, the flexible distribution gives the smallest deviations at the quantiles examined, of the same order as the variability estimated by resampling, but its support excludes 0.95\% of the days, precisely the calmest ones. For the temperature, the deviations of the normal distribution are smaller than the sampling variability of the quantiles, and the two-parameter model remains a reasonable choice for this description. The examples show that a good fit requires examining the support and the goodness of fit, not only the moments, and that the choice of model depends on the physical question. All results are reproduced by a notebook provided with the paper, and a suggested classroom sequence, with tasks and expected answers, is included.

\medskip
\noindent\textbf{Keywords}: generalized lambda distribution; quantile function; method of moments; meteorological data; statistics education; Python.
\end{abstract}
\end{otherlanguage}
